\section{Discussion and Limitations}
\label{sec:discussion}

\paragraph{The causal ablation isolates a mechanism, not a full explanation.}
The noise-fill comparison in \cref{subsec:e-ablation} shows that texture
preservation, specifically, drives the advantage on no-tissue-interface
structures; it does not by itself explain the smaller, more uniform gains
observed on tissue-interface structures, which likely stem from the wider
space of region-to-region contrast relationships (including the sign
inversions in \cref{eq:affine}) that no single real acquisition protocol
supplies. Disentangling these two effects further is left to future work.

\paragraph{Why the dissociation follows from the anatomy.} Preserving real
intensity structure matters where the target is defined by tissue
\emph{appearance}, and not where it is defined by an anatomical
\emph{interface}. Liver, spleen and kidney are encapsulated organs bounded by
fat planes -- a genuine physical interface, which is why edge-following methods
such as geodesic active contours have long succeeded on
them~\cite{kavur2021chaos,lee2007gaclivers}; an augmentation that scrambles
internal texture while preserving that interface loses little. Diffuse glioma
has no such interface: tumour cells infiltrate beyond both the
contrast-enhancing margin and the peritumoral FLAIR
abnormality~\cite{gliomamargins2025,flairectomy2022}, so the target must be
inferred from appearance rather than followed along an edge. MS lesions are
likewise defined by tissue signal characteristics rather than an anatomical
border~\cite{zhang2008mstexture}. This account is an interpretation of the
ablation, not an independent proof of it; the causal evidence is the
intervention in \cref{subsec:e-ablation}. It also extends to the ON-Harmony
result once its actual labels are examined (\cref{subsec:e-ablation}): its
31 regions -- ventricles, cortex, white matter, deep grey nuclei -- are
bounded by genuine tissue-type transitions throughout, so despite first
appearances it groups with the organ tasks rather than the no-interface
ones. The larger negative magnitude there, compared to CHAOS's near-null
one, is plausibly a density effect -- near-total brain volume relabelled at
once, rather than a handful of organs -- which we leave to a dedicated
follow-up with a second densely-labelled task.

\paragraph{The label step needs only the annotation training already
assumes.} PALETTE's one label-consuming step uses whatever segmentation labels
the training set happens to carry -- sparse, partial, or task-specific -- and
never a dense anatomical parcellation. This is the practical difference from
label-generative synthesis: SynthSeg's dense requirement means the augmentation
costs an annotation richer than the task, whereas here the step is free at
training time and can be switched off entirely on an unlabelled dataset, in
which case PALETTE degrades gracefully to its unsupervised steps.

\paragraph{Scope.} PALETTE targets the single-source setting: one labelled
training modality, no access to target-domain data of any kind. It is
input-level and architecture-agnostic, and we do not claim it subsumes
multi-source domain generalization methods that have access to several
labelled source domains; the two are complementary rather than competing
solutions to the same problem.

\section{Conclusion}
\label{sec:conclusion}

We asked when preserving real image texture matters for contrast-robust
segmentation, and found that the answer is a property of the task rather than
of the augmentation: it matters where the target is defined by tissue
appearance, and not where it is defined by an anatomical interface. A
controlled intervention -- holding the partition and per-region statistics
fixed and swapping only label-style noise fill for a real-intensity remap --
separates our no-tissue-interface and tissue-interface tasks along exactly
that line, worth $+7.7$ and $+7.2$ Dice on the former and nothing ($+1.1$,
$-1.1$) on the latter -- a pattern a third, far more densely-labelled
interface task (healthy-brain) also follows, just with a larger negative
swing ($-4.5$, $-1.4$) plausibly reflecting how much more of the volume is
relabelled at once (\cref{subsec:e-ablation}). PALETTE, the single-source
augmentation built on this observation, partitions the foreground by
unsupervised $k$-means and a Voronoi sub-parcellation and applies a signed
affine remap of the real intensities within each sub-region, requiring
neither a dense anatomical label map nor any target-domain data. Across
eight cross-contrast and cross-modality settings plus a fifth,
single-modality liver-tumour task, and cross-dataset external test sets,
spanning abdominal, brain, brain-tumor, and liver imaging, a single nnU-Net
configuration trained with PALETTE attains the best Dice in every task and
improves significantly over all five competing methods on Dice, remaining
at or near the best on boundary distance without always being the outright
aggregate winner there. We regard the
task-level account as the more durable contribution: it predicts in advance
where a texture-preserving augmentation will and will not pay off, and that
prediction is testable on any new task.

%%%%%%%%% ACKNOWLEDGEMENTS - if applicable, uncomment for camera-ready
% \paragraph{Acknowledgements.} This work was enabled in part by support
% provided by the Digital Research Alliance of Canada (alliancecan.ca) and
% computational resources on the Killarney/Vulcan clusters.
